\section{Who checks the monitor, intervention and composition?}\label{sec:verify}
A green dashboard can mean that no release occurred, or that the release did not reach the dashboard. Checking the dashboard's repeatability cannot distinguish these cases. The relevant object is the \emph{joint} law of the observation and the true event. For a pure deterministic monitor to reveal an event exactly, every physical trajectory delivering the same monitor record must carry the same event value on the represented scope. A single conflicting pair refutes exact decoding. For uncertain monitors, retain all joint laws compatible with the available calibration and event semantics and optimize the desired error over them (Theorems~\ref{thm:t25} and~\ref{ext:t25j}). Counts from the monitor alone identify only its marginal law.

The engineer might install a second logger to check the first. That intervention may change timing, permissions or the path taken by the job. To relate its result to the original operation one needs a comparison of physical histories \emph{under both procedures}. If the starts and successive histories can be coupled, given the matched past, so that the event still agrees until the first discrepancy, then the probability of a discrepancy bounds the difference in event probabilities. Conditional row checks can supply this calculation. Merely comparing frequencies after the logger was installed cannot reconstruct the counterfactual (Theorems~\ref{thm:t26} and~\ref{ext:t26c}).

Finally, one can validate a build component and a release component in isolation and still miss a violation at their handoff. A certificate for the composition needs the permitted joint calls, their record meanings and the physical relation between local and global events. As a numerical illustration, if two local events each have probability \(1/10\), their union can have probability anywhere from \(1/10\) to \(1/5\) without a dependence assumption. A global event that is not their union requires a different event relation altogether. Theorems~\ref{thm:t27} and~\ref{ext:t27j} compute the range in the stipulated joint model. These proofs direct verification effort to the boundary where one component's output becomes another's permission.
